% !TEX root = print.tex

% Master blueprint document for the s-numbers project.
%
% This file is intended for compilation with the `leanblueprint`
% tool (https://github.com/PatrickMassot/leanblueprint), which
% turns the LaTeX source plus a Lean 4 project into a hyperlinked
% web blueprint that tracks the formalisation progress.

\documentclass[11pt]{report}

\usepackage[utf8]{inputenc}
\usepackage{amsmath, amssymb, amsthm}
\usepackage{xcolor}
\usepackage[colorlinks,
            linkcolor=blue!55!black,
            citecolor=blue!55!black,
            urlcolor=blue!55!black]{hyperref}
\usepackage[capitalize]{cleveref}
\usepackage{etoolbox}

% --- leanblueprint macros -------------------------------------------------
% Definitions for compiling with pdflatex; in the web build the `blueprint`
% package shipped with leanblueprint provides the real, clickable versions.
\providecommand{\lean}[1]{}
\providecommand{\leanok}{}
\providecommand{\uses}[1]{}
\providecommand{\notready}{}
\providecommand{\proves}[1]{}

% \lean{A, B}: name the Lean declaration(s) a statement corresponds to. The
% tag sits at the top of the environment (the leanblueprint convention), but
% reads better after the statement, so it is stored here and printed by
% \bp@leanflush, which \AtEndEnvironment hooks to the theorem environments.
% Two subtleties in the printing:
%   * `\detokenize` is required: declaration names contain underscores, which
%     TeX would otherwise read as math subscripts;
%   * the list is split at the commas, so that a long list wraps while each
%     individual name stays on one line (detokenized spaces are not
%     breakpoints, so a single detokenized string would run into the margin).
\makeatletter
\newcommand{\bp@leanlist}{}
\renewcommand{\lean}[1]{\gdef\bp@leanlist{#1}}
\newcommand{\bp@leanname}[1]{\mbox{\detokenize{#1}}}
\newcommand{\bp@leanflush}{%
  \ifx\bp@leanlist\@empty\else
    \par\nobreak\smallskip
    {\normalfont\footnotesize\raggedright
     [\textsf{Lean:}\nobreakspace\ttfamily
     \expandafter\@for\expandafter\bp@leanitem\expandafter:%
       \expandafter=\bp@leanlist\do{%
         \expandafter\bp@leanname\expandafter{\bp@leanitem}\allowbreak}%
     \unskip\unpenalty]%
     \par}%
    \medskip
    \gdef\bp@leanlist{}%
  \fi}
\makeatother

% \leanok marks a formalized statement or proof. Everything in this blueprint
% is formalized, so printing a marker on every one of them would be noise;
% the convention is stated once in the first chapter instead.
\renewcommand{\leanok}{}

% Give TeX extra inter-word stretch before it lets a line overflow into the
% margin; long typewriter identifiers and inline formulas make this necessary.
\emergencystretch=3em

% --- theorem environments -------------------------------------------------
\newtheorem{theorem}{Theorem}[chapter]
\newtheorem{proposition}[theorem]{Proposition}
\newtheorem{lemma}[theorem]{Lemma}
\newtheorem{corollary}[theorem]{Corollary}
\newtheorem{remark}[theorem]{Remark}
\theoremstyle{definition}
\newtheorem{definition}[theorem]{Definition}
\newtheorem{example}[theorem]{Example}

% Print the Lean declaration at the end of each statement. The hook runs
% while the theorem's italic font is still in force, hence the \normalfont
% inside \bp@leanflush.
\makeatletter
\AtEndEnvironment{theorem}{\bp@leanflush}
\AtEndEnvironment{proposition}{\bp@leanflush}
\AtEndEnvironment{lemma}{\bp@leanflush}
\AtEndEnvironment{corollary}{\bp@leanflush}
\AtEndEnvironment{definition}{\bp@leanflush}
\makeatother

% --- title ----------------------------------------------------------------
\title{Blueprint:\\ s-Numbers of Bounded Linear Operators\\
       between Banach Spaces}
\author{Mario Ullrich}
\date{\today}

\begin{document}
\maketitle
\tableofcontents

% !TEX root = print.tex

\chapter{Setting and prerequisites}
\label{ch:setting}

Throughout the blueprint, $\mathbb{K}$ denotes a non-trivially normed
field — in practice $\mathbb{R}$ or $\mathbb{C}$ — and $X, Y, Z, W$ are
Banach spaces over $\mathbb{K}$. We write $\mathcal{L}(X, Y)$ for the
space of bounded linear operators $X \to Y$ with the operator norm.

\begin{definition}[Rank of a continuous linear map]
\label{def:rank}
\lean{ContinuousLinearMap.rank}
\leanok
For $S \in \mathcal{L}(X, Y)$ the \emph{rank} of $S$ is the cardinal
dimension of its range,
\[
\operatorname{rank}(S) := \dim_{\mathbb{K}} \operatorname{range}(S)
                        \in \mathrm{Card}.
\]
\end{definition}

\chapter{Axioms of an $s$-number sequence}
\label{ch:axioms}

\begin{definition}[$s$-number sequence; Pietsch axioms]
\label{def:s-number-sequence}
\lean{SNumbers.IsSNumberSequence}
\leanok
\uses{def:rank}
A family
\[
s \;=\; (s_n)_{n \in \mathbb{N}}, \qquad
s_n : \mathcal{L}(X, Y) \longrightarrow [0, \infty),
\]
indexed for all Banach pairs $(X, Y)$ over $\mathbb{K}$ is called an
\emph{$s$-number sequence} if:
\begin{description}
  \item[(S1)]\label{ax:S1} $\|S\| = s_0(S) \ge s_1(S) \ge \dots \ge 0$
        for every $S$.
  \item[(S2)]\label{ax:S2} $s_n(S + T) \le s_n(S) + \|T\|$.
  \item[(S3)]\label{ax:S3} $s_n(BSA) \le \|B\| \, s_n(S) \, \|A\|$ for
        composable bounded operators $A, B$.
  \item[(S4)]\label{ax:S4} $s_n(S) = 0$ whenever $\operatorname{rank}(S) \le n$.
  \item[(S5)]\label{ax:S5} $s_n(\mathrm{id}_{\ell_2^{n+1}}) = 1$.
\end{description}
The bundled Lean structure splits (S1) into the three predicates
\texttt{norm\_at\_zero}, \texttt{antitone}, \texttt{nonneg}.
\end{definition}

\begin{definition}[Strict $s$-number sequence]
\label{def:strict-s-number-sequence}
\lean{SNumbers.IsStrictSNumberSequence}
\leanok
\uses{def:s-number-sequence}
An $s$-number sequence is \emph{strict} if it additionally satisfies the
strengthening
\begin{description}
  \item[(S5$'$)] $s_n(\mathrm{id}_X) = 1$ for every Banach space $X$ with
        $\dim X > n$, not merely on $X = \ell_2^{n+1}$.
\end{description}
\end{definition}

\chapter{The five canonical examples}
\label{ch:examples}

We construct the five textbook $s$-number sequences. Each definition is
encoded as a \texttt{noncomputable def} returning a real number; the proof
that it satisfies the axioms (S1)–(S5) is the content of an
\texttt{IsSNumberSequence} theorem.

\section{Approximation numbers $a_n$}
\label{sec:approx}

\begin{definition}[Approximation number]
\label{def:approximation-number}
\lean{SNumbers.approximationNumber}
\leanok
\uses{def:rank}
For $S \in \mathcal{L}(X, Y)$ and $n \in \mathbb{N}$,
\[
a_n(S) := \inf \bigl\{ \|S - A\| : A \in \mathcal{L}(X, Y),
                       \operatorname{rank}(A) \le n \bigr\}.
\]
\end{definition}

\begin{theorem}[Approximation numbers form a strict $s$-number sequence]
\label{thm:approx-is-snumber}
\lean{SNumbers.isStrictSNumberSequence_approximationNumber}
\leanok
\uses{def:approximation-number, def:strict-s-number-sequence}
The family $a = (a_n)_n$ satisfies all axioms (S1)--(S5) and the
strengthening (S5$'$): $a_n(\mathrm{id}_X) = 1$ on every Banach space
$X$ with $\dim X > n$.
\end{theorem}
\begin{proof}
Axioms (S1)--(S4) are proved in
\texttt{SNumbers/Approximation.lean}: \texttt{approximationNumber\_nonneg},
\texttt{approximationNumber\_antitone}, \texttt{approximationNumber\_le\_norm},
\texttt{approximationNumber\_add\_le},
\texttt{approximationNumber\_comp\_comp\_le} (using
\texttt{rank\_comp\_comp\_le} and \texttt{Real.sInf\_smul\_of\_nonneg}),
and \texttt{approximationNumber\_eq\_zero\_of\_rank\_le}. The strict
normalisation (S5$'$) is \texttt{approximationNumber\_strict}:
for any $L$ with $\operatorname{rank}(L) \le n$, the range of $L$ is a
proper finite-dimensional, hence closed, subspace of $X$; Riesz's lemma
provides $x_0 \notin \operatorname{range}(L)$ with $r \|x_0\| \le \|x_0 - L
x_0\|$ for any $r < 1$, so $\|\mathrm{id} - L\| \ge \|(\mathrm{id}-L) x_0\|
/ \|x_0\| \ge r$; letting $r \to 1$ gives $\|\mathrm{id} - L\| \ge 1$,
and combined with the trivial upper bound $a_n(\mathrm{id}) \le \|\mathrm{id}\| = 1$
we obtain equality. The classical (S5) on $\ell_2^{n+1}$ specialises (S5$'$).
The argument needs neither SVD nor Auerbach; the blueprint stubs in
\texttt{BasicResults.SVD} are kept for the future identification
$a_n(S) = \sigma_{n+1}(S)$ for compact $S$.
\end{proof}

\section{Bernstein numbers $b_n$}
\label{sec:bern}

\begin{definition}[Bernstein number]
\label{def:bernstein-number}
\lean{SNumbers.bernsteinNumber}
\leanok
For $S \in \mathcal{L}(X, Y)$ and $n \in \mathbb{N}$,
\[
b_n(S) := \sup_{\dim M = n+1}
            \inf_{\substack{x \in M \\ \|x\| = 1}} \|Sx\|.
\]
\end{definition}

\begin{proposition}[$b_n$ is an $s$-number sequence]
\label{prop:bernstein-is-snumber}
\lean{SNumbers.isSNumberSequence_bernsteinNumber}
\notready
\uses{def:bernstein-number, def:s-number-sequence}
\end{proposition}

\section{Gelfand numbers $c_n$}
\label{sec:gelf}

\begin{definition}[Gelfand number]
\label{def:gelfand-number}
\lean{SNumbers.gelfandNumber}
\leanok
For $S \in \mathcal{L}(X, Y)$ and $n \in \mathbb{N}$,
\[
c_n(S) := \inf \bigl\{ \|S \restriction_M\| :
            M \subseteq X \text{ closed subspace, } \operatorname{codim} M \le n
            \bigr\}.
\]
\end{definition}

\begin{proposition}[$c_n$ is an $s$-number sequence]
\label{prop:gelfand-is-snumber}
\lean{SNumbers.isSNumberSequence_gelfandNumber}
\notready
\uses{def:gelfand-number, def:s-number-sequence}
\end{proposition}

\section{Kolmogorov numbers $d_n$}
\label{sec:kol}

\begin{definition}[Kolmogorov number]
\label{def:kolmogorov-number}
\lean{SNumbers.kolmogorovNumber}
\leanok
For $S \in \mathcal{L}(X, Y)$ and $n \in \mathbb{N}$,
\[
d_n(S) := \inf_{\substack{V \subseteq Y \\ \dim V \le n}}
            \sup_{\|x\| \le 1} \inf_{v \in V} \|Sx - v\|.
\]
\end{definition}

\begin{theorem}[$d_n$ is a strict $s$-number sequence]
\label{thm:kolmogorov-is-snumber}
\lean{SNumbers.isStrictSNumberSequence_kolmogorovNumber}
\leanok
\uses{def:kolmogorov-number, def:strict-s-number-sequence}
The Kolmogorov numbers satisfy all five Pietsch axioms (S1)–(S5) and
moreover the strengthening (S5$'$): $d_n(\mathrm{id}_X) = 1$ on every
Banach space $X$ with $\dim X > n$. The argument requires the scalar
field $\mathbb{K}$ to be a \texttt{DenselyNormedField} (in particular
$\mathbb{R}$ or $\mathbb{C}$) and additionally \texttt{CompleteSpace}
$\mathbb{K}$ for the (S5$'$) step (which closes finite-dimensional
subspaces via Riesz's lemma).
\end{theorem}

\section{Hilbert numbers $h_n$}
\label{sec:hil}

For the Hilbert numbers we work over $\mathbb{R}$ and use
$L_2 := \ell_2(\mathbb{N}; \mathbb{R})$.

\begin{definition}[Hilbert number]
\label{def:hilbert-number}
\lean{SNumbers.hilbertNumber}
\leanok
For $S \in \mathcal{L}(X, Y)$ and $n \in \mathbb{N}$,
\[
h_n(S) := \sup \bigl\{ a_n(B \, S \, A) :
            A : L_2 \to X, \ B : Y \to L_2, \ \|A\|, \|B\| \le 1 \bigr\}.
\]
\end{definition}

\begin{proposition}[$h_n$ is an $s$-number sequence]
\label{prop:hilbert-is-snumber}
\lean{SNumbers.isSNumberSequence_hilbertNumber}
\notready
\uses{def:hilbert-number, def:approximation-number, def:s-number-sequence}
\end{proposition}

\chapter{Pietsch's sandwich theorem}
\label{ch:sandwich}

\begin{theorem}[Hilbert-space coincidence]
\label{thm:all-eq-on-hilbert}
\lean{SNumbers.allSNumbers_eq_on_HilbertSpace}
\notready
\uses{def:s-number-sequence, def:approximation-number}
On a Hilbert space, every $s$-number sequence agrees with the
approximation numbers: $s_n(S) = a_n(S)$ for every continuous linear map
$S$ between Hilbert spaces.
\end{theorem}

\begin{theorem}[Pietsch \emph{sandwich}]
\label{thm:sandwich}
\lean{SNumbers.hilbertNumber_le_sn_le_approximationNumber}
\uses{thm:approx-is-snumber, prop:hilbert-is-snumber}
For every $s$-number sequence $s$ on real Banach spaces and every operator
$S$,
\[
h_n(S) \;\le\; s_n(S) \;\le\; a_n(S).
\]
On Hilbert spaces all three numbers agree.
\end{theorem}
\begin{proof}
The upper bound $s_n(S) \le a_n(S)$ is proved without \texttt{sorry}
in \texttt{SNumbers.sn\_le\_approximationNumber}: for any $A$ with
$\operatorname{rank}(A) \le n$ axiom (S2) yields
$s_n(S) \le s_n(A) + \|S - A\| = \|S - A\|$ by (S4); take the infimum.
The lower bound $h_n(S) \le s_n(S)$ uses (S3) plus the Hilbert-space
coincidence theorem and is currently \texttt{sorry}.
\end{proof}

\chapter{Basic auxiliary results}
\label{ch:basic-results}

The companion library \texttt{BasicResults} (sibling of \texttt{SNumbers})
collects two generic functional-analysis results that are independent of
the s-numbers framework. The first — Auerbach's lemma — is fully proved;
the second — compact-operator SVD — is a blueprint stub.

\begin{theorem}[Auerbach's lemma]
\label{thm:auerbach}
\lean{exists_isAuerbachBasis}
\leanok
Every finite-dimensional normed space over $\mathbb{R}$ of positive
dimension admits a basis $(e_i)$ with $\|e_i\| = 1$ whose dual coordinate
functionals also satisfy $\|e_i^*\| = 1$.
\end{theorem}

\begin{theorem}[SVD for compact operators]
\label{thm:compact-svd}
\lean{SVD.compactOperator_svd}
\notready
For every compact operator $S : H_1 \to H_2$ between Hilbert spaces
there exist orthonormal sequences $(u_k) \subseteq H_1$,
$(v_k) \subseteq H_2$ and a non-increasing sequence of non-negative reals
$\sigma_1 \ge \sigma_2 \ge \dots \ge 0$ with
$S x = \sum_k \sigma_k \langle x, u_k \rangle v_k$ for every $x \in H_1$.
\end{theorem}


\end{document}
